\documentclass[10pt,a4paper]{amsart}
\usepackage[margin=19mm]{geometry}
\usepackage{amsmath,amssymb,mathtools}
\usepackage[scr=rsfs,cal=euler]{mathalfa}
\usepackage{xfrac}
\usepackage[unicode,hidelinks,bookmarks=false]{hyperref}
\newcommand{\ie}{i.e.\ }
\newcommand{\cf}{cf.\ }
\newcommand{\resp}{respectively\ }
\newcommand{\Subsection}{\subsection}
\providecommand{\nobreakdash}{\nobreak\hbox{-}}
\setlength{\emergencystretch}{2em}
\multlinegap0pt
\usepackage{bm}
\usepackage{bbm}
\usepackage{dsfont}
\usepackage{xspace}
\usepackage{cancel}
\usepackage{mathtools}
\usepackage{amsthm}
\makeatletter
\@ifundefined{lemma}{\newtheorem{lemma}{Lemma}}{}
\makeatother
\DeclarePairedDelimiter{\abs}{\lvert}{\rvert} % Absolute value. \abs* does auto-sizing, \abs[\big] etc does fixed sizing
\newcommand{\BB}{\mathcal{B}}
\renewcommand{\AA}{\mathcal{A}}
\renewcommand{\geq}{\geqslant}
\renewcommand{\leq}{\leqslant}
\title[Non-Homotopic Drawings of Multigraphs]{Non-Homotopic Drawings of Multigraphs}
\author{Ant\'onio Gir\~ao, Freddie Illingworth, Alex Scott, David R. Wood}
\date{Source: arXiv:2401.10615v2 (CC BY 4.0)}
\begin{document}
\maketitle

\noindent\emph{Source and licence.} Non-Homotopic Drawings of Multigraphs. Version arXiv:2401.10615v2 (CC BY 4.0), distributed under the Creative Commons Attribution 4.0 International licence, as recorded in the arXiv licence metadata for this exact version and shown on the abstract page https://arxiv.org/abs/2401.10615v2. Full rights notice: licenses/arxiv-2401.10615-CC-BY-4.0.txt.\par\smallskip

\noindent\textit{Editorial scope.} Counted unit: the Lemma whose source label is \texttt{sequencebound} in the article (source0.tex lines 510--515, source label \texttt{sequencebound}), delivered together with the article's own complete original proof of it (source0.tex lines 517--554, one \texttt{proof} environment of 38 lines). No mathematics is written, repaired or re-derived: statement and proof are the article's own text line for line. Required context retained uncounted: none. Every definition, display and label of those lines is retained unchanged. Omitted from this package and not redistributed: 1--509, 516--516, 555--778. Nothing inside a retained excerpt is omitted or altered: every retained body is delivered exactly as the article's lines are written, so no omission receipt of any kind accompanies this package. Citations: the retained excerpts carry no citation command at all, so no bibliography entry of the article is transcribed and no reference is redistributed. Reference adaptation: none was needed and none was made; every reference inside the delivered unit resolves inside it, so no unresolved reference or citation is produced. Printed slips kept literal: no printed word, symbol, hypothesis or number of the counted statement or of its proof was altered. Layout and class: the article's own class is \texttt{article} with packages including geometry, parskip, wrapfig, fontenc, inputenc, amsfonts, amsmath, amssymb, mathtools, microtype, anyfontsize, mathpazo, xcolor, url; the wrapper is the lane's pinned \texttt{amsart} editorial wrapper at 10pt on A4, which restates the theorem-like environments the retained text needs and adds the article's own macro definitions that the retained text needs (\texttt{\textbackslash AA}, \texttt{\textbackslash BB}, \texttt{\textbackslash abs}, \texttt{\textbackslash geq}, \texttt{\textbackslash leq}), with extra wrapper packages (bm, bbm, dsfont, xspace, cancel, mathtools). The delivered numbering is the unit's own. Assets: no retained span contains a figure environment, a graphics-inclusion command, a PDF-page inclusion, a table or a TikZ picture; no image file is copied into this package, the package declares no asset dependency and no image file is redistributed with it. Licence: version arXiv:2401.10615v2 of this article is distributed under the Creative Commons Attribution 4.0 International licence, as recorded in the arXiv licence metadata for this exact version and shown on the abstract page https://arxiv.org/abs/2401.10615v2. A full rights notice is delivered at licenses/arxiv-2401.10615-CC-BY-4.0.txt. Characters: every retained excerpt is pure ASCII, so the delivered engine, XeLaTeX with its default Latin Modern text font, reports no missing character.
\begin{lemma}\label{sequencebound}
    Let $k < n$ be positive integers. Let $\AA$ be a set of different monotone drawing sequences of length $n$, each pair of which cross at most $k$ times. Then
    \begin{equation*}
        \abs{\AA} \leq 2 \cdot \binom{2n}{k + 1}.
    \end{equation*}
\end{lemma}

\begin{proof}
    Let $g(n, k)$ denote the size of the largest set of length $n$ monotone drawing sequences where every pair crosses at most $k$ times. In this definition we allow $k \geq n$, in which case $g(n, k) = g(n, n - 1)$, since no two monotone drawing sequences of length $n$ cross $n$ times. We also allow $k = -1$, in which case $g(n, -1) = 1$ (as no two sequences cross at most $-1$ times).

    Let $\AA$ be a set of monotone drawing sequences of length $n + 1$, such that every pair crosses at most $k$ times. The proof will proceed by deleting the first entry of each sequence in $\AA$ and analysing how many sequences elide.

    Let $\BB$ be those sequences of $\AA$ which, upon deleting their first entry, are no longer monotone drawing sequences. The only sequence in $\BB$ is $0, \ast, \ast, \dots, \ast$ and so $\abs{\BB} \leq 1$. Let $\AA_{+}, \AA_{0}, \AA_{-}, \AA_{\ast}$ be those sequences in $\AA - \BB$ which start with a $+$, $0$, $-$, $\ast$, respectively. Let $\AA'_{+}$ be those sequences obtained by deleting the first entry (which, of course, is a $+$) from sequences in $\AA_{+}$. Similarly define $\AA'_{0}, \AA'_{-}, \AA'_{\ast}$. Note that each of these is a set of monotone drawing sequences of length $n$. These sets may not be disjoint. Let $\AA' = \AA'_{+} \cup \AA'_{0} \cup \AA'_{-} \cup \AA'_{\ast}$ be those sequences obtained by deleting the first entry from a sequence in $\AA - \BB$. Note $\AA'$ is a set of monotone drawing sequences of length $n$ where every pair of sequences crosses at most $k$ times and so $\abs{\AA'} \leq g(n, k)$.

    In any monotone drawing sequence, the symbol $\ast$ can be only be succeeded or preceded by a $0$ or $\ast$. In particular, the first entry of a sequence in $\AA'_{\ast}$ is a $\ast$ or a $0$. If a sequence in $\AA'_{+} \cup \AA'_{0} \cup \AA'_{-}$ starts with a $\ast$, then the original sequence (before deleting the first entry) must have been $0, \ast, \ast, \dotsc, \ast$ and this was discounted when we removed $\BB$. Further, the only sequence in $\AA'_{+} \cup \AA'_{0} \cup \AA'_{-}$ starting with a $0$ is $0, \ast, \ast, \dotsc, \ast$. Hence, $\abs{\AA'_{\ast} \cap (\AA'_{+} \cup \AA'_{0} \cup \AA'_{-})} \leq 1$.

    We now consider those sequences that are in $\AA'_{+} \cap (\AA'_{0} \cup \AA'_{-})$ and show that every pair of them crosses at most $k - 1$ times. Indeed, suppose that $(a_i), (b_i) \in \AA'_{+} \cap (\AA'_{0} \cup \AA'_{-})$ cross $k$ times. Let $j$ be the first index where $a_j \neq b_j$.  We first show that neither $a_j$ nor $b_j$ is a $\ast$. Suppose that $a_j = \ast$. Since $(a_i) \in \AA'_{+} \cup \AA'_{0} \cup \AA'_{-}$, its first entry is a $+$, $0$, or $-$. Hence, from the definition of monotone drawing sequences, there is some $\ell < j$ with $a_\ell = 0$. By the minimality of $j$, $b_\ell$ is also $0$. But $(b_i) \in \AA'_{+} \cup \AA'_{0} \cup \AA'_{-}$, so every entry after $b_\ell$ is a $\ast$ and so $b_j = \ast$, contradicting $a_j \neq b_j$. 

    Since $a_j$, $b_j$ are distinct and neither is $\ast$, $(a_i)$ is either above or below $(b_i)$ at $j$. Without loss of generality $(a_i)$ is below $(b_i)$ at $j$. But $(a_i)$ and $(b_i)$ agree before $j$ and cross $k$ times, so there are indices $j = i_1 < i_2 < \dotsb < i_{k + 1}$ with $(a_i)$ below $(b_i)$ at $i_1, i_3, \dotsc$ and $(a_i)$ above $(b_i)$ at $i_2, i_4, \dotsc$. As $(a_i) \in \AA'_{+}$, the sequence obtained by appending a $+$ to the start of $(a_i)$ is in $\AA$, while, since $(b_i) \in \AA'_{0} \cup \AA'_{-}$, the sequence obtained by appending a $0$ or $-$ (depending on whether $(b_i) \in \AA'_{0}$ or $(b_i) \in \AA'_{-}$) to the start of $(b_i)$ is in $\AA$. These two sequences cross $k + 1$ times, a contradiction.

    Hence, $\AA'_{+} \cap (\AA'_{0} \cup \AA'_{-})$ is a set of monotone drawing sequences of length $n$, every pair of which cross at most $k - 1$ times and so $\abs{\AA'_{+} \cap (\AA'_{0} \cup \AA'_{-})} \leq g(n, k - 1)$. Similarly, $\abs{\AA'_{-} \cap (\AA'_{0} \cup \AA'_{+})} \leq g(n, k - 1)$. Hence, the number of sequences in at least two of $\AA'_{+}, \AA'_{0}, \AA'_{-}, \AA'_{\ast}$ is at most $2 g(n, k - 1) + 1$. By the inclusion-exclusion principle we have
    \begin{align*}
       \abs{\AA}  
       = \abs{\AA'_{+}} + \abs{\AA'_{0}} + \abs{\AA'_{-}} + \abs{\AA'_{\ast}} + \abs{\BB} 
        & \leq \abs{\AA'} + 2g(n, k - 1) + 2 \\
        & \leq g(n, k) + 2g(n, k - 1) + 2.
    \end{align*}
    Hence, $g(n + 1, k) \leq g(n, k) + 2g(n, k - 1) + 2$. It suffices to show that $g(n, k) \leq 2 \cdot \binom{2n}{k + 1}$ for all $-1 \leq k \leq n - 1$. There are some easy base cases:
    \begin{itemize}
        \item $g(n, -1) = 1$ for all $n$ (any two sequences cross at least 0 times);
        \item $g(1, k) = 1$ for all $k \geq 0$ (the only monotone drawing sequence of length 1 is $0$);
        \item $g(2, k) = 4$ for all $k \geq 0$ (the only monotone drawing sequences of length 2 are $+, 0$; $0,0$; $-, 0$; and $0,\ast$).
    \end{itemize}
    These and the recurrence give $g(n, 0) \leq 4n - 4$ for $n \geq 2$. Hence, the required inequality holds whenever $n \leq 2$ or $k \leq 0$. Suppose the result holds for $n \geq 2$. First suppose that $1 \leq k \leq n - 1$. Then
    \begin{align*}
        g(n + 1, k) & \leq g(n, k) + 2g(n, k - 1) + 2 \leq 2 \tbinom{2n}{k + 1} + 4 \tbinom{2n}{k} + 2 \\
        & \leq 2 \tbinom{2n}{k + 1} + 4 \tbinom{2n}{k} + 2 \tbinom{2n}{k - 1} = 2 \tbinom{2n + 2}{k + 1}.
    \end{align*}
    Finally suppose that $k = n$. Now
    \begin{align*}
        g(n + 1, n) & \leq g(n, n) + 2g(n, n - 1) + 2 = 3g(n, n - 1) + 2 \\
        & \leq 6 \tbinom{2n}{n} + 2 \leq 2 \tbinom{2n + 2}{n + 1},
    \end{align*}
    where the last inequality holds as $\binom{2n + 2}{n + 1} \binom{2n}{n}^{-1} = (4n + 2)/(n + 1) > 3$ for $n \geq 2$.
\end{proof}

\end{document}
